\chapter{Tinjauan Pustaka}

\section{Landasan Teori}
Bagian ini memuat teori yang menjadi dasar penelitian. Konsep dan sifat formal
dapat dinyatakan dengan lingkungan definisi, teorema, dan bukti berbahasa Indonesia.

\begin{definisi}[Proxy Means Test]
Misalkan $y_i$ menyatakan log pengeluaran per kapita rumah tangga ke-$i$ dan
$\mathbf{x}_i$ vektor karakteristik teramati. \emph{Proxy Means Test} adalah fungsi
prediksi $\hat{f}:\mathbf{x}_i \mapsto \hat{y}_i$ yang diperoleh dengan meminimumkan
suatu fungsi galat $\sum_i L\!\left(y_i,\hat{f}(\mathbf{x}_i)\right)$.
\end{definisi}

\begin{teorema}\label{teo:ols}
Pada model linear $y = \mathbf{X}\beta + \varepsilon$ dengan
$\mathrm{E}(\varepsilon \mid \mathbf{X}) = \mathbf{0}$, penduga kuadrat terkecil
$\hat{\beta} = (\mathbf{X}^\top\mathbf{X})^{-1}\mathbf{X}^\top y$ bersifat takbias.
\end{teorema}

\begin{proof}
Substitusi $y = \mathbf{X}\beta + \varepsilon$ memberikan
$\hat{\beta} = \beta + (\mathbf{X}^\top\mathbf{X})^{-1}\mathbf{X}^\top\varepsilon$.
Karena $\mathrm{E}(\varepsilon \mid \mathbf{X}) = \mathbf{0}$, maka
$\mathrm{E}(\hat{\beta} \mid \mathbf{X}) = \beta$.
\end{proof}

\section{Penelitian Terkait}
Beberapa penelitian terdahulu mengembangkan PMT dengan pendekatan
\emph{machine learning} \parencite{etang2022,taufiq2021}.

\section{Penulisan Kutipan}
Format sitasi mengikuti APA versi 7 dengan penyesuaian bahasa Indonesia
(\emph{dan} untuk dua pengarang, \emph{dkk.} untuk lebih dari dua pengarang,
\emph{hal.} untuk halaman).

\subsection{Kutipan Tidak Langsung (Parafrase)}
Kutipan tidak langsung mencantumkan nama pengarang dan tahun \emph{tanpa} nomor
halaman. Satu pengarang: \textcite{cochran1991} memformalkan teori penarikan
sampel. Dua pengarang: \textcite{taufiq2021} membandingkan beberapa algoritma.
Lebih dari dua pengarang: \textcite{etang2022} menekankan pentingnya validasi
silang. Versi dalam kurung berturut-turut: \parencite{cochran1991};
\parencite{taufiq2021}; \parencite{etang2022}.

\subsection{Kutipan Langsung Pendek (kurang dari atau sama dengan 40 kata)}
Kutipan langsung pendek diberi tanda kutip serta nomor halaman. Naratif (nomor
halaman di akhir kutipan): \textcite{cochran1991} menyatakan bahwa ``penarikan
sampel acak sederhana memberikan estimasi takbias bagi rata-rata populasi''
(hal.~10). Dalam kurung: ``replikasi data memungkinkan akses yang konsisten
lintas waktu'' \parencite[hal.~781]{taufiq2021}.

\subsection{Kutipan Langsung Panjang (lebih dari 40 kata)}
Kutipan langsung panjang ditulis pada paragraf tersendiri, menjorok 1 tab dari
batas kiri teks, tanpa tanda kutip:
\begin{kutipanpanjang}
Proxy means testing merupakan pendekatan yang banyak dipakai untuk menargetkan
bantuan sosial ketika data pendapatan rumah tangga tidak tersedia atau sulit
diverifikasi; metode ini menaksir kesejahteraan dari sejumlah peubah proksi yang
mudah diamati, sehingga keputusan penargetan dapat dilakukan secara lebih objektif
dan dapat direplikasi \parencite[hal.~25]{etang2022}.
\end{kutipanpanjang}

\subsection{Mengutip dari Kutipan (Sitasi Sekunder)}
Diperbolehkan bila pustaka asli sulit ditemukan: nama pengarang asli dicantumkan
pada kalimat, sumber sekunder (tempat kutipan ditemukan) pada akhir kalimat.
Menurut Nelder dan Wedderburn (1972), seperti dikutip dalam
\textcite{cochran1991}, prinsip estimasi takbias menjadi dasar banyak metode
inferensi. Pada Daftar Pustaka hanya dicantumkan sumber sekunder yang
benar-benar dibaca (Cochran), bukan sumber asli (Nelder dan Wedderburn).

\subsection{Sitasi Lembaga/Instansi Pemerintah}
Bila penulis berupa lembaga/instansi, pada \textbf{kemunculan pertama} ditulis
nama lengkap diikuti akronim, lalu \textbf{kemunculan selanjutnya} cukup akronim
(pedoman 4.4). Pada entri bib: \verb|author = {{Badan Pusat Statistik}}| (kurung
ganda agar tidak dipenggal) dan \verb|shortauthor = {BPS}|.

Kemunculan pertama (naratif): \textcite{bps2024} menetapkan garis kemiskinan
secara nasional. Dalam kurung: \parencite{bps2024}. Kemunculan selanjutnya cukup
akronim: \textcite{bps2024}, atau \parencite{bps2024}. Pada Daftar Pustaka, nama
lembaga \emph{selalu} ditulis lengkap (Badan Pusat Statistik), bukan akronim.

\subsection{Sitasi Direktorat di Bawah Kementerian}
Bila karya disusun oleh suatu \textbf{direktorat di bawah kementerian}, ikuti
APA-7: unit \textbf{paling spesifik menjadi penulis}, kementerian induk menjadi
\textbf{penerbit} (jangan menggabung keduanya sebagai satu nama). Pada entri bib:
\verb|author = {{Direktorat Jenderal Pembangunan Desa dan Perdesaan}}|,
\verb|shortauthor = {Ditjen PDP}|, dan \verb|institution = {Kementerian ...}|.

Kemunculan pertama: \textcite{ditjenpdp2023} menyusun pedoman alokasi dana desa.
Selanjutnya: \textcite{ditjenpdp2023}; atau dalam kurung \parencite{ditjenpdp2023}.
Pada Daftar Pustaka, nama direktorat (penulis) diikuti judul, lalu kementerian
sebagai penerbit.
